
\begin{result}{B.22}{Proved}{A geometric extension through the square--triangle families}
There is a specified realization of the five object families
\[
 \mathrm I_{10},\quad\mathrm I_6,\quad\mathrm I_4,
 \quad\mathrm {II}_2,\quad\mathrm {II}'_2
\]
and the six arrow families
\[
 f_{B,s},\quad f_{\vec Q,a},\quad\tau_{\pi,s},
 \quad\gamma_{\pi,s},\quad\gamma''_{\pi,s}
\]
(the first row of arrows has two families) with the following
counts and complete presentation:
\[
\begin{array}{c|r}
 \text{objects}&12+20+30+60+60=182\\
 \text{arrows}&60+120+60+60+60=360\\
 \text{free cycles after a spanning tree}&179\\
 \text{defining relators}&180.
\end{array}                                                     \tag{B.45}
\]
The comparison sends these paths to $\Pi_1M_{0,5}$ and is an
equivalence after imposing the specified relators.  The
quarter-turn and the two erasures are furnished by the
one-node plumbing family below.  We specify $f_{\vec Q,a}$
as an erasure followed by a core path; this records its
path convention, in addition to its endpoint incidence.
All $360$ braid images and all $180$ relators are computed
in \texttt{b\_first\_tier.json}.
\end{result}
\begin{proof}
Use
\[
 V(q)=(\infty,-q,q,-1,1),\qquad
 q\in\mathbb C\setminus\{0,1,-1\},\qquad
 \epsilon=1/10.                                      \tag{B.46}
\]
At $q=0$ this has a stable limit consisting of the square
four-point component $(\infty,-1,1,0)$ and a three-point
tail.  Indeed, on the tail put $w=z/q$, so its marks are
$-1,1,\infty$, and the node coordinates $z,v=1/w$ satisfy
$zv=q$.  Thus $q$ is the actual normal plumbing coordinate.
For $|q|\leq\epsilon$ one may choose disjoint coordinate
discs of radius $1/2$ at the two branches.  The points
$V(\epsilon)$ and $V(i\epsilon)$ give smoothing
representatives of the right and left square--triangle
objects.  Both have the stabilizer $H$ in (B.33), so each
has sixty labelled realizations.

The twist is the positive quarter-turn
$q(t)=\epsilon\exp(\pi it/2)$; in the braid computation it
is replaced by
$\epsilon(1-t+it)$.  The radial interpolation between these
two paths avoids $0,1,-1$, proving the required endpoint
homotopy before normalization.  The two erasures are the
positive real and imaginary intervals
\[
\begin{split}
 \gamma_{\pi,s}:&\quad
  q:\epsilon\longrightarrow\sqrt5-2,\\
 \gamma''_{\pi,s}:&\quad
  q:i\epsilon\longrightarrow i/\sqrt3.
\end{split}                                                     \tag{B.47}
\]
They begin in the indicated normal ray, then continue
in the displayed smooth family.  The imaginary endpoint
is the bipyramid of B.20.  The real endpoint is pentagonal.
To check this last assertion, the transformation
$(z+1)/(z-1)$ sends the fifth roots of unity to $\infty$
and the four imaginary points of magnitudes
$\cot(\pi/5)$ and $\cot(2\pi/5)$.  Their ratio is
\[
 \frac{\cot(\pi/5)}{\cot(2\pi/5)}=2+\sqrt5.
\]
After scaling the larger magnitude to $1$ the smaller
is $\sqrt5-2$, giving (B.47).  Its cyclic order on the
real projective line is $(0,3,1,2,4)$, up to reversal,
and its stabilizer is
\[
 H_P=\langle(0\,3\,1\,2\,4),(12)(34)\rangle\simeq D_5.

overall\tag{B.48}
\]
Consequently the real erasure has the exact endpoint map
$\mathfrak S_5/H\to\mathfrak S_5/H_P$, with five incoming
arrows per pentagon.  The imaginary erasure has
$\mathfrak S_5/H\to\mathfrak S_5/H_B$, with three incoming
arrows per bipyramid.  The twist is the identity on the
common coset parameter $\mathfrak S_5/H$.

For clarity, the $120$ choices denoted $f_{\vec Q,a}$
are specified as follows.  A left square--triangle object
is the pair $(B,P)$, with $P$ one of the three equatorial
marks of its bipyramid after erasure.  Choose one of the
other two equatorial marks $s$, and use the composite
\[
 (B,P)\xrightarrow{\gamma''}B\xrightarrow{f_{B,s}}Q.
                                                               \tag{B.49}
\]
The apex of $Q$ is $s$.  The remaining equatorial pair
of $B$ is a diagonal of its oriented square.  Starting
at $P$ and following the cyclic orientation gives a side
$a$ of $Q$.  Conversely a side $a$ of an oriented square,
with initial vertex $P$, determines that diagonal, one
of the two incident bipyramids, and the pair $(B,P)$.
This is a bijection between (B.49) and the parameters
$(\vec Q,a)$: two choices per source and four per target.
It proves the finite endpoint comparison as well as
specifying the path used.  The source does not give a
separate parameterization of its direct $f_{\vec Q,a}$
path; (B.49) is the explicit convention of this realization.

Here is the complete relator prescription, including its
finite indexing and its computed word data.  Keep the first
fifty objects and sixty arrows of B.21.  Append the twelve
pentagons, sixty right square--triangles and sixty left
square--triangles, in lexicographic coset order.  Append
the sixty real erasures, sixty twists, sixty imaginary
erasures, then the $120$ composites (B.49), in that order.
The latter are ordered by their source index and then by
the core-edge index.  This numbers every object and arrow
without a choice.

Retain B.21's forty-nine-edge tree.  In increasing order,
add a new edge to the tree precisely when it joins two
previously different components.  The result has $181$
edges and all $182$ vertices.  Write $t_v$ for its unique
rooted path.  For each non-tree edge $E$, let
\[
 g_E=t_{s(E)}E t_{t(E)}^{-1}
\]
in traversal order, and compute its word $W_E(a,b,x,y,z)$
by the normalized-path and combing algorithms of B.21.
For the $168$ new non-tree edges impose exactly
\[
 g_E= W_E(A,B,X,Y,Z),                                \tag{B.50}
\]
where (B.38) gives the words on the right in the original
eleven core cycles.  Add the twelve core relators
(B.39)--(B.40).  These are the $180$ relators in (B.45).
The file \texttt{b\_first\_tier.json} lists every $E$,
$W_E$, tree path, and both sides of each expanded relation;
none of the $168$ words in (B.50) is left as an uncomputed
input.

For the path calculation, the normalized coordinates are
again (B.41)--(B.42), with $r=-iq$.  The real erasures
have no projected crossings, since all four finite
coordinates remain real and distinct.  For each pentagon
choose its transport through the first incoming real
erasure and its already specified right-object transport.
The right and left object transports use B.21's same
least-integer detour rule.  The sixty twists and sixty
imaginary erasures have rational coordinate paths;
all $120$ are homotopic to one chord by the strict
half-plane tests recorded in the program.  Exact Sturm
counts, exact crossing signs over $\mathbb Q(\sqrt3)$,
and Artin combing verify their words.  Pentagon coordinates
in $\mathbb Q(\sqrt5)$ are checked algebraically for the
shared endpoints; they introduce no crossings or numerical
approximations in the braid calculation.  Each composite
(B.49) is then checked by multiplying its two verified
words.

Finally, collapse the chosen tree.  The resulting free
group has $360-182+1=179$ generators: eleven old and
$168$ new.  Relations (B.50) eliminate all the new
ones.  The remaining presentation is exactly B.21,
which is $\Gamma_{0,5}$.  The comparison and its inverse
on vertex groups are therefore the actual path functor
and (B.38).  Tree transports extend the inverse to the
whole groupoid.  This proves both the relator prescription
and its completeness.
\end{proof}

\begin{result}{B.23}{Proved}{The realized paths come from an explicit four-point double-cover functor}
Put
\[
 U=\mathbb P^1_q\setminus\{0,1,-1,\infty\},\qquad
 p:U\longrightarrow M_{0,4},\quad p(q)=q^2,
\]
using the coordinate $M_{0,4}=\mathbb P^1_\lambda
\setminus\{0,1,\infty\}$.  There are exact functors
\[
 \Pi_1U\xrightarrow{p_*}\Pi_1M_{0,4},\qquad
 \Pi_1U\xrightarrow{D_*}\Pi_1M_{0,5},\quad
 D(q)=(\infty,-q,q,-1,1).                         \tag{B.51}
\]
The first is the two-sheeted covering functor.  The second
is induced by the double cover of the four-marked sphere
branched at $0,\infty$, retaining the point above $\infty$
and the four ordered points above $1,\lambda$.
For every arrow of the realized $360$-arrow model, its
image is a composite of explicitly relabelled images
under $D_*$ of lifted four-point paths.  Thus this is a
proved realization through a family of relative dimension
one, with the covering and forgetting maps retained.
It does not identify an unspecified sewing-only convention
for the other arrows in the margin of scan 33.

For a positive real basepoint $0<q<1$, order the four
finite points as $-1,-q,q,1$.  If $u_0,u_+,u_-$ are the
standard meridians at $q=0,1,-1$ specified below, their
braid images are
\[
 u_0\longmapsto\sigma_2^2,\qquad
 u_+\longmapsto\sigma_1^2\sigma_3^2,\qquad
 u_-\longmapsto\sigma_2\sigma_1^2\sigma_3^2\sigma_2^{-1}.
                                                               \tag{B.52}
\]
Under the scalar quotient and B.19's generators these are
\[
 b,\qquad b^{-1}a^{-1}z^{-1}y^{-1}x^{-1}z,
 \qquad az^{-1}yz,                                  \tag{B.53}
\]
respectively.
\end{result}
\begin{proof}
The map $q\mapsto q^2$ is unramified on the displayed
open curves and has two points above each $\lambda$.
The universal four-point sphere is normalized to
$(0,\infty,1,\lambda)$.  Its double cover $z^2=w$
has branch points $0,\infty$, and the four chosen lifts
are $\pm1,\pm q$.  Keep the ramification point $\infty$
as the fifth mark and omit the other ramification point.
The square-root choice $q$ and the chosen lift $+1$ give
the exact labels in (B.51).  This constructs the two
maps of spaces and therefore their path-groupoid functors.

Let $\ell_0,\ell_1$ be meridians of $0,1$ in the
$\lambda$-line, based at $q^2$.  Covering monodromy is
$\ell_0\mapsto1$ and $\ell_1\mapsto0$ in $\mathbb Z/2$.
The two-vertex lifted graph of the standard two-loop
spine has the three free generators
\[
 u_0=\ell_0^2,\qquad u_+=\ell_1,\qquad
 u_-=\ell_0\ell_1\ell_0^{-1}.
\]
Take the stems to $q=1$ along the positive real interval.
The stem to $q=-1$ first turns $q$ positively through a
half-circle around $0$ and then follows the negative
real interval.  The loop around $q=0$ rotates the inner
two points once, giving $\sigma_2^2$.  A loop around
$q=1$ simultaneously turns the two disjoint adjacent
pairs at $-1$ and $1$, giving
$\sigma_1^2\sigma_3^2$.  The chosen half-circle stem
interchanges the inner two points by $\sigma_2$; this
gives the last word in (B.52).  These observations can
also be checked directly from the four trajectories
$-1,-q,q,1$; the disjoint small meridian discs make the
two simultaneous twists commute.

The middle braid has linking number $1$ on strands
$1,2$ and the others have linking number zero.  Removing
that one full scalar braid and applying the exact combing
algorithm yields (B.53).  The program
\texttt{b\_cover\_map.py} carries out these three word
calculations and verifies the braid identities.

For the core arrows the lifted parameter path is
$q=ir$, $1/\sqrt3\leq r\leq1$, whose projection is
the real interval $\lambda=-r^2$.
For the one-node families it is exactly (B.46)--(B.47),
and for the quarter-turn it is
$q=\epsilon e^{\pi it/2}$, projecting to
$\lambda=\epsilon^2e^{\pi it}$.  Formula (B.49)
supplies the remaining composites.  This gives the
claimed expression for every arrow in the realized
subdiagram by actual four-point covering paths,
with no dimension-only identification.
\end{proof}
